% Exact CI-verified theorem; Frontier Continuation run 33345071604.
\subsection{A seven-sector bicolour algebra on the $216$ five-circuits}
The maximal-overlap relation between the $216$ sentinel five-circuits and the $540$ sentinel six-circuits,
\[
|C_5\cap C_6|=3,
\]
is $(20,8)$-biregular and splits into two $PSp(4,3)$ orbitals of size $2160$.  Let their $(216\times540)$ biadjacency matrices be $M_+$ and $M_-$.  Each colour is $(10,4)$-biregular, and exact integer computation gives
\[
M_+M_+^T=M_-M_-^T=10I+A_{30},
\]
\[
M_+M_-^T+M_-M_+^T=4A_{20},
\]
where $A_{30}$ and $A_{20}$ are simple regular graphs of degrees $30$ and $20$ on the five-circuit set.  Their edge supports are disjoint and
\[
[A_{30},A_{20}]=0.
\]
Therefore the uncoloured incidence matrix $M=M_++M_-$ satisfies
\[
MM^T=20I+2A_{30}+4A_{20}.
\]
The exact spectra are
\[
\operatorname{Spec}(A_{30})=30^1+12^{15}+6^{39}+0^{60}+(-4)^{81}+(-6)^{20},
\]
\[
\operatorname{Spec}(A_{20})=20^1+8^{24}+2^{80}+(-2)^{81}+(-4)^{15}+(-10)^{15}.
\]
The separating operator $A_{30}+7A_{20}$ has seven distinct eigenvalues and forces seven simultaneous sectors, with $(A_{30},A_{20};\dim)$ equal to
\[
(-4,-2;81),\ (0,2;60),\ (6,8;24),\ (-6,2;20),\ (12,-10;15),\ (6,-4;15),\ (30,20;1).
\]
Hence the joint spectral dimensions are
\[
\boxed{1,15,15,20,24,60,81},
\]
and the maximal-overlap Gram spectrum is
\[
\boxed{\operatorname{Spec}(MM^T)=160^1+64^{24}+28^{60}+16^{35}+4^{96}}.
\]
In particular $M$ has full row rank $216$ over characteristic zero.  The proof is exact: integer annihilating polynomials and trace multiplicities replace floating-point diagonalisation.  These seven spaces are certified $PSp(4,3)$-invariant spectral sectors of the incidence algebra; identification with irreducible complex representations is deliberately left to a separate character-theoretic audit.
